% ===================================================================
%  TABLICA Nr 14 — Ulica. — Kupcy.
%  Meme regime que les precedents. Tableau a UNE SEULE COLONNE, et le
%  plus peuple du livret : quatre-vingt-dix-neuf renvois, presque
%  tous des metiers.
%
%  QUATRE MOTS-OUTILS SONT EN GRAS DANS LE FAC-SIMILE : « Til » au 1,
%  « la » et « lu » au 13, « e, » au 16. Le polonais n'a pas
%  d'article ; le tcheque non plus, et il deplace le gras sur le mot
%  plein voisin — « \VUgras{vojenské} hudby » pour « \VUgras{la}
%  militistal muziko », « \VUgras{ten} » pour « \VUgras{lu} ». Cette
%  colonne fait pareil. On ne supprime pas un gras du fac-simile, on
%  lui trouve un porteur.
%
%  ET « tintisto » (53) EST LE TEINTURIER, NON L'ETAMEUR. Le piege
%  est reel : le tcheque a lu « stano », l'etain, et a ecrit
%  « cinar ». L'ido dit « tinto », la teinture — farbiarz. Le
%  neerlandais dit verver, et la planche montre des cuves.
%
%  Trois autres faux amis du meme tableau : « mastro-veturi » (57)
%  sont les voitures de maitre, non les vehicules de demenagement ;
%  « regulierino » (89) une religieuse ; « kompana siorino » (95) une
%  dame de compagnie, non une compagne.
% ===================================================================

%%K t14-apar-1 apar
\VUsaut{4.0mm}
\VUtitre{15.0pt}{60}{TABLICA Nr 14}
\VUsaut{3.0mm}

%%K t14-tit-1 sub
\VUpk{11.4pt}{40}{Ulica. --- Kupcy.}
\VUsaut{3.0mm}

%%K t14-01-1 p
1. --- Mój przyjaciel Jan, który mieszka na wsi, przyjechał spędzić trzy dni w mojej
rodzinie. \VUgras{Do} tej pory nie miał sposobności zobaczyć wielkiego miasta; dlatego
spędzamy wszystkie nasze dni na spacerach, a ja tłumaczę mu to, czego nie zna.

%%K t14-02-1 p
2. --- Najpierw poszliśmy zwiedzić \VUgras{obserwatorium} \textsuperscript{(1)}, które
bardzo go zajęło. Wracając \VUgras{ulicą} \textsuperscript{(3)}, gdzie jest
\VUgras{łaźnia turecka} \textsuperscript{(2)}, spotkaliśmy \VUgras{pogrzeb}
\textsuperscript{(4)}. Przed \VUgras{konduktem żałobnym} szło \VUgras{duchowieństwo,}
poprzedzając \VUgras{karawan,} w którym umieszczono \VUgras{trumnę.} Wielu przyjaciół
towarzyszyło rodzinie \VUgras{w żałobie.}

%%K t14-02-2 p
Po przejściu orszaku przeszliśmy przez ulicę przed wielką \VUgras{piwiarnią}
\textsuperscript{(5)}, w której wielu \VUgras{konsumentów} piło \VUgras{piwo} na
tarasie, osłoniętych oszkloną \VUgras{markizą} \textsuperscript{(6)}.

%%K t14-03-1 p
3. --- Jan był bardzo zdziwiony \VUgras{tarczą herbową} umieszczoną między sklepem
\VUgras{likiernika} \textsuperscript{(7)} a sklepem \VUgras{handlarza nutami}
\textsuperscript{(10)}. Zapytał mnie o jej znaczenie; odpowiedziałem, że wskazuje ona
angielski \VUgras{konsulat} \textsuperscript{(8)}, i zwróciłem jego uwagę na
\VUgras{konsula} \textsuperscript{(9)}, który jest na balkonie.

%%K t14-tit-2 sub
\VUpk{10.2pt}{40}{Zwiedzanie sklepów.}
\VUsaut{3.0mm}

%%K t14-04-1 p
4. --- Oczywiście zatrzymaliśmy się przed wystawą \VUgras{instrumentów muzycznych,}
\VUgras{fisharmonii,} \VUgras{organów} i \VUgras{fortepianów}; oglądaliśmy ilustrowane
okładki \VUgras{partytur} i \VUgras{utworów muzycznych} na fortepian i podziwialiśmy
\VUgras{lirę} ze złoconego drewna, która służy za szyld.

%%K t14-05-1 p
5. --- Tumany kurzu robione przez \VUgras{zamiataczy} \textsuperscript{(11)} i
\VUgras{wóz zamiatający} \textsuperscript{(12)} zmusiły nas, żeby szybko przejść przed
pięknymi \VUgras{kompletami mebli} \VUgras{handlarza mebli} \textsuperscript{(13)} i
przed wystawą \VUgras{piecyków} i \VUgras{lamp} \VUgras{blacharza}
\textsuperscript{(14)}.

%%K t14-06-1 p
6. --- Zresztą w tym miejscu byliśmy zmuszeni zejść z chodnika, bo był on zatarasowany
paczkami, które brano z \VUgras{wozu przeprowadzkowego} \textsuperscript{(16)}, żeby
umieścić je w \VUgras{magazynie} \textsuperscript{(15)}, który właśnie wynajęto.

%%K t14-06-2 p
To przeszkodziło nam zobaczyć piękne powozy \VUgras{karetnika} \textsuperscript{(17)},
ale nie przeszkodziło nam zatrzymać się, żeby popatrzeć na \VUgras{pakowacza}
\textsuperscript{(19)}, który robił skrzynię dla swojego sąsiada, \VUgras{handlarza
rowerów i samochodów} \textsuperscript{(18)}. Potem, ponieważ było już nieco późno,
wróciliśmy do domu.

%%K t14-tit-3 sub
\VUpk{10.2pt}{40}{Przechadzka po mieście.}
\VUsaut{3.0mm}

%%K t14-07-1 p
7. --- Nazajutrz moja matka zaproponowała Janowi, żeby go sfotografowano, potem poleciła
mi, żebym towarzyszył mu do \VUgras{fotografa} \textsuperscript{(20)}. Po obiedzie
weszliśmy do \VUgras{pracowni} artysty, gdzie musieliśmy dość długo czekać. Żeby się
czymś zająć, oglądaliśmy \VUgras{portrety} \textsuperscript{(22)} wiszące na ścianie.

%%K t14-07-2 p
Kiedy przyszła nasza kolej, praca nie trwała długo, a zaraz gdy mój przyjaciel skończył
pozować przed \VUgras{aparatem fotograficznym} \textsuperscript{(21)}, zeszliśmy.

%%K t14-08-1 p
8. --- Na pierwszym piętrze zobaczyliśmy przez drzwi \VUgras{fryzjera}
\textsuperscript{(23)}, które są otwarte, jego czeladnika, który strzygł włosy klienta.

%%K t14-08-2 p
Wyszedłszy na ulicę, minęliśmy \VUgras{sklep tytoniowy} \textsuperscript{(24)}. Jan był
tak rozbawiony czerwoną \VUgras{latarnią} \textsuperscript{(25)} i \VUgras{grubą fajką},
która służy za \VUgras{szyld} \textsuperscript{(26)}, że zderzył się z dwoma starymi
panami, którzy rozmawiali, \VUgras{zażywając tabaki} \textsuperscript{(27)}.

%%K t14-tit-4 sub
\VUpk{10.2pt}{40}{Cukiernia.}
\VUsaut{3.0mm}

%%K t14-09-1 p
9. --- \VUgras{Banknoty} i \VUgras{monety} ze wszystkich krajów wystawione w
\VUgras{witrynie} \VUgras{wekslarza} \textsuperscript{(29)} przyciągnęły na kilka chwil
naszą uwagę, ale ponieważ była godzina podwieczorku, weszliśmy do \VUgras{cukiernika}
\textsuperscript{(31)}, nie zatrzymując się dłużej.

%%K t14-10-1 p
10. --- Widząc wszystkie \VUgras{łakocie,} które mieścił sklep, \VUgras{ciasta,
pierniki, biszkopty} i wszelkiego rodzaju \VUgras{cukierki,} nie mogliśmy łatwo wybrać.
W końcu Jan wybrał \VUgras{ciastka z kremem,} a ja wziąłem tylko małą \VUgras{tartę
wiśniową,} bo od słodyczy \VUgras{bolą mnie zęby,} a wcale nie życzę sobie chodzić za
często do \VUgras{dentysty} \textsuperscript{(30)}.

%%K t14-tit-5 sub
\VUpk{10.2pt}{40}{Pisanie na maszynie.}
\VUsaut{3.0mm}

%%K t14-11-1 p
11. --- Żeby pokazać mojemu przyjacielowi \VUgras{maszynę do pisania}, o której słyszał,
ale której nie znał, weszliśmy do \VUgras{biura} \textsuperscript{(32)} mojego
\VUgras{kuzyna} \textsuperscript{(34)}. Zastaliśmy go dyktującego swoje listy swojemu
\VUgras{sekretarzowi} \textsuperscript{(35)}, który jest \VUgras{stenografem}. Ledwie
list był skończony, \VUgras{maszynista} \textsuperscript{(33)} przepisywał go na swojej
maszynie. Nie chcąc przerywać zajęć mojego krewnego, zostaliśmy u niego tylko przez
kilka chwil.

%%K t14-12-1 p
12. --- Pod kantorem mojego kuzyna mieszka wielki \VUgras{zegarmistrz-jubiler}
\textsuperscript{(37)}, który jest ponadto złotnikiem. Jego wspaniały sklep mieści wiele
\VUgras{zegarów, zegarów wahadłowych, zegarków kieszonkowych,} \VUgras{łańcuszków} i
cennych \VUgras{klejnotów,} na przykład \VUgras{naszyjników perłowych,} złotych
\VUgras{bransolet, kolczyków} i \VUgras{pierścionków} zdobionych \VUgras{kamieniami
szlachetnymi} i \VUgras{diamentami}. Jedna z witryn jest osobno zarezerwowana dla
\VUgras{wyrobów złotych} i mieści ważny zbiór srebrnych \VUgras{zastaw stołowych}.

%%K t14-13-1 p
13. --- Skręcając na rogu ulicy, zauważyłem, że zmieniono \VUgras{tabliczkę}
\textsuperscript{(36)} umieszczoną blisko \VUgras{numeru} domu. Miałem właśnie powiedzieć
głośno swoją uwagę, kiedy zabrzmiały wesołe dźwięki \VUgras{wojskowej} muzyki.
Pobiegliśmy naprzeciw pułkowi, nie zastanawiając się, że \VUgras{ten} ma przejść przed
sklepami \VUgras{kapelusznika} \textsuperscript{(39)} i \VUgras{rękawicznika}
\textsuperscript{(40)}, i że bylibyśmy równie dobrze ustawieni, żeby zobaczyć jego
defiladę, gdybyśmy zostali przed witryną \VUgras{optyka} \textsuperscript{(38)},
całkiem zaopatrzoną w \VUgras{binokle, lorgnony} i \VUgras{lunetki}.

%%K t14-tit-6 sub
\VUpk{10.2pt}{40}{Defilada marszowa.}
\VUsaut{3.0mm}

%%K t14-14-1 p
14. --- Przed \VUgras{pułkiem} \textsuperscript{(41)} szedł \VUgras{tamburmajor}
\textsuperscript{(42)}, za którym szli \VUgras{dobosze} \textsuperscript{(43)},
\VUgras{trębacze} \textsuperscript{(44)} i cała \VUgras{orkiestra}.
\VUgras{Generał} \textsuperscript{(45)}, który był na placu ze swoim adiutantem,
zatrzymał się blisko \VUgras{wodotrysku} \textsuperscript{(49)} \VUgras{fontanny}
\textsuperscript{(48)}, żeby popatrzeć na \VUgras{defiladę}.

%%K t14-14-2 p
\VUgras{Oficerowie sztabowi} \textsuperscript{(46)}, \VUgras{pułkownik} i
\VUgras{podpułkownik} pozdrowili go szpadą. \VUgras{Majorowie} byli przed swoimi
\VUgras{batalionami} \textsuperscript{(47)}, a każdy \VUgras{kapitan} dowodził swoją
\VUgras{kompanią,} podczas gdy \VUgras{porucznicy,} \VUgras{podporucznicy} i
\VUgras{podoficerowie} zajmowali swoje zwykłe miejsce przy swoich ludziach.

%%K t14-15-1 p
15. --- Wielu przechodniów zebrało się na \VUgras{skrzyżowaniu} \textsuperscript{(50)};
kupcy, \VUgras{parasolnik} \textsuperscript{(51)}, \VUgras{farbiarz}
\textsuperscript{(53)}, \VUgras{rusznikarz} \textsuperscript{(54)} wyszli, żeby zobaczyć
\VUgras{żołnierzy}. Wszystkie pojazdy, \VUgras{dorożki} \textsuperscript{(52)},
\VUgras{powozy pańskie} \textsuperscript{(57)}, a także \VUgras{beczkowóz}
\textsuperscript{(56)}, ustawiły się wzdłuż chodnika.

%%K t14-tit-7 sub
\VUpk{10.2pt}{40}{Sceny na ulicy.}
\VUsaut{3.0mm}

%%K t14-15-2 p
\VUgras{Komiwojażer} \textsuperscript{(55)}, niosący w ręku swoje \VUgras{kufry z
próbkami,} szybko przeszedł przez ulicę, a osoby siedzące na \VUgras{imperiale}
\textsuperscript{(59)} \VUgras{omnibusu} \textsuperscript{(58)} odwróciły się, żeby
nacieszyć się widowiskiem. Biedny \VUgras{śpiewak uliczny} \textsuperscript{(60)} ze
swoją \VUgras{katarynką} \textsuperscript{(61)} musiał przerwać swoją
\VUgras{piosenkę,} bo huk bębnów zagłuszył jego głos.

%%K t14-16-1 p
16. --- Kiedy pułk przybył przed \VUgras{sklep z tkaninami} \textsuperscript{(64)},
\VUgras{szwaczki} \textsuperscript{(63)} i \VUgras{krawcy} \textsuperscript{(62)}
zostawili swoje \VUgras{maszyny do szycia} i swoją pracę, żeby patrzeć przez okno.
\VUgras{Kupiec} \textsuperscript{(65)}, który właśnie ustawił nowe \VUgras{ubrania}
\textsuperscript{(66)} na swojej \VUgras{wystawie,} musiał kazać wnieść do środka
\VUgras{sukna} \textsuperscript{(67)} i \VUgras{płótna} ustawione na chodniku, blisko
\VUgras{otworu kanału} \textsuperscript{(68)}, w obawie, że tłum je przewróci; nawet
stary \VUgras{domokrążca} \textsuperscript{(69)}, od którego dozorczyni targowała
pończochy, był zmuszony wejść do sieni, żeby dokończyć swoją sprzedaż.

%%K t14-16-2 p
Towarzyszyliśmy pułkowi aż do koszar \VUgras{i,} bardzo zadowoleni ze swojego dnia,
wróciliśmy o godzinie kolacji.

%%K t14-tit-8 sub
\VUpk{10.2pt}{40}{Rozmaite rzemiosła i zawody.}
\VUsaut{3.0mm}

%%K t14-17-1 p
17. --- Następnego dnia mój ojciec zaproponował nam zwiedzenie drukarni tutejszej
gazety. Przyjęliśmy to z zapałem.

%%K t14-17-2 p
\VUgras{Drukarz} jest też \VUgras{księgarzem} \textsuperscript{(70)}
\VUgras{(=sprzedawcą książek),} \VUgras{wydawcą,} \VUgras{handlarzem papieru} i
\VUgras{introligatorem}. Umieścił na pierwszym piętrze \VUgras{redakcję}
\textsuperscript{(71)} gazety, a także \VUgras{pracownię rytowniczą,} w której pracują
\VUgras{litografowie} \textsuperscript{(72)} i \VUgras{miedziorytnicy}
\textsuperscript{(73)}, wreszcie \VUgras{zecernię,} w której są \VUgras{drukarze}
\textsuperscript{(74)}. Właściwa \VUgras{drukarnia} \textsuperscript{(75)},
\VUgras{prasy drukarskie} i rozmaite maszyny są na parterze.

%%K t14-tit-9 sub
\VUpk{10.2pt}{40}{Jan zadaje bardzo wiele pytań.}
\VUsaut{3.0mm}

%%K t14-18-1 p
18. --- Wychodząc z drukarni, Jan bardzo zdziwiony zatrzymał się przed małą oszkloną
skrzynką umieszczoną na słupie, naprzeciw \VUgras{sklepu pasmanteryjnego}
\textsuperscript{(77)}, i zapytał, do czego ona służy. Mój ojciec wytłumaczył mu, że to
\VUgras{ostrzegacz pożarowy} \textsuperscript{(76)}. Zresztą wszystko w mieście było dla
mojego przyjaciela zdumieniem, od zwykłej \VUgras{kamiennej fontanny}
\textsuperscript{(78)} i lekkiego \VUgras{trójkołowca dostawczego}
\textsuperscript{(80)} prowadzonego przez \VUgras{boya} \textsuperscript{(79)} w liberii,
aż po \VUgras{wóz pocztowy} \textsuperscript{(81)}.

%%K t14-19-1 p
19. --- Podczas gdy mój ojciec zostawił nas na kilka chwil, żeby porozmawiać z
\VUgras{poborcą podatkowym} \textsuperscript{(83)}, który zatrzymał się, żeby
porozmawiać z \VUgras{adwokatem} \textsuperscript{(82)} i \VUgras{notariuszem}
\textsuperscript{(84)}, bawiliśmy się, śledząc oczami ruch na ulicy. Najrozmaitsze osoby
przechodziły przed nami: \VUgras{czeladnik cukierniczy} \textsuperscript{(85)}, niosący
w koszu wspaniały \VUgras{pasztet,} szedł za \VUgras{handlarzem odzieży}
\textsuperscript{(86)}; \VUgras{latarnik} \textsuperscript{(87)} rozmawiał z
\VUgras{gazownikiem} \textsuperscript{(88)}, \VUgras{zakonnica} \textsuperscript{(89)},
trzymająca za rękę sierotę, wracała do swojego klasztoru.

%%K t14-tit-10 sub
\VUpk{10.2pt}{40}{W drodze powrotnej do domu.}
\VUsaut{3.0mm}

%%K t14-20-1 p
20. --- W chwili, gdy mój ojciec dołączył do nas, Jan roześmiał się głośno, widząc
brudną twarz i dziwaczne ubranie dwóch \VUgras{kominiarzy} \textsuperscript{(91)}
całkiem osmolonych.

%%K t14-21-1 p
21. --- Chciałem kupić u \VUgras{kwiaciarki} \textsuperscript{(92)} roślinę dla mojej
matki, ale mój ojciec zwrócił mi uwagę, że byłaby bardzo ciężka do niesienia i że lepiej
byłoby nabyć \VUgras{pomarańcze}. Zbliżyliśmy się do \VUgras{sprzedawczyni}
\textsuperscript{(94)}, a kiedy \VUgras{wychowawczyni} \textsuperscript{(93)}, która
wybierała dla swojego wychowanka, skończyła swój zakup, kazaliśmy się obsłużyć.

%%K t14-22-1 p
22. --- Wracając do domu, spotkaliśmy kilku znajomych: starą przyjaciółkę mojej
rodziny, która opierała ramię na swojej \VUgras{damie do towarzystwa}
\textsuperscript{(95)}, \VUgras{inżyniera} \textsuperscript{(96)} od mostów i dróg, i
jedną z moich kuzynek z jej \VUgras{nianią} \textsuperscript{(98)}.

%%K t14-22-2 p
Potem spotkaliśmy \VUgras{modystkę} \textsuperscript{(97)} mojej matki i dwóch
\VUgras{mnichów} \textsuperscript{(99)} obcych miastu, którzy podeszli do nas, żeby
zapytać mojego ojca o drogę na dworzec.
\VUsaut{6.0mm}
\VUfilet{22mm}
